\documentclass[12pt]{article}
\usepackage[utf8]{inputenc}
\usepackage[margin=1in]{geometry}
\usepackage{amsmath,amssymb,amsthm,amsfonts,mathrsfs}
\usepackage{hyperref}

\title{The Eight-Fold Path to a Classical P\'olya Urn Computation}
\author{Christopher D. Long}
\date{}

\newtheorem{theorem}{Theorem}
\newtheorem{lemma}{Lemma}
\newtheorem{remark}{Remark}

\begin{document}

\maketitle

\begin{abstract}
We consider the classical two-color P\'olya urn started from $3$ white balls and $2$ black balls. At each step, one ball is drawn uniformly at random, returned to the urn, and an additional ball of the same color is added. If $X_N$ denotes the proportion of white balls after $N$ draws, we compute
\[
P_\infty:=\lim_{N\to\infty}\mathbb{P}\!\left(X_N>\frac12\right).
\]
The answer is
\[
P_\infty=\frac{11}{16}.
\]
This note presents eight rigorous viewpoints on the computation: an exact beta-binomial formula, a Riemann-sum limit, a Bayesian/conjugacy argument, a martingale-plus-factorial-moments proof, an exchangeability/de Finetti proof, a continuous-time branching-process embedding, a first-tie/Catalan decomposition, and an analytic generating-function approach based on a transport PDE together with a Beta-kernel normalization.
\end{abstract}

\tableofcontents

\section{Introduction and notation}

Let $W_0=3$ and $B_0=2$. After $N$ draws, let
\[
W_N=\text{number of white balls}, \qquad B_N=\text{number of black balls},
\]
so that
\[
W_N+B_N=N+5.
\]
Let
\[
I_N:=W_N-3
\]
be the number of white draws among the first $N$ draws. Then
\[
W_N=3+I_N,\qquad B_N=2+(N-I_N).
\]
The white proportion is
\[
X_N:=\frac{W_N}{N+5}=\frac{3+I_N}{N+5}.
\]
We seek
\[
P_\infty:=\lim_{N\to\infty}\mathbb{P}\!\left(X_N>\frac12\right).
\]

Note that
\[
X_N>\frac12
\quad\Longleftrightarrow\quad
\frac{3+I_N}{N+5}>\frac12
\quad\Longleftrightarrow\quad
I_N>\frac{N-1}{2}.
\]

For $a>0$ and $m\in\mathbb{N}$, write
\[
(a)_m:=a(a+1)\cdots(a+m-1),
\]
with $(a)_0=1$, for the rising factorial.

The eight routes below fall into three broad families:
\begin{itemize}
    \item exact discrete formulas (Solutions 1, 7, 8),
    \item probabilistic representation theorems and stochastic-process arguments (Solutions 3, 4, 5, 6),
    \item asymptotic extraction (Solution 2).
\end{itemize}
Several of them identify the full limit law
\[
X_\infty\sim \mathrm{Beta}(3,2),
\]
from which the desired probability is the Beta tail
\[
\int_{1/2}^1 12x^2(1-x)\,dx=\frac{11}{16}.
\]

\begin{theorem}
For the P\'olya urn started from $(3,2)$,
\[
\lim_{N\to\infty}\mathbb{P}\!\left(X_N>\frac12\right)=\frac{11}{16}.
\]
\end{theorem}

\section{Solution 1: Exact state probabilities (the beta-binomial law)}

This gives the exact finite-$N$ distribution.

\subsection*{Step 1: probability of a specific color word}
Fix a particular color sequence of length $N$ containing exactly $k$ white draws and $N-k$ black draws.

Along that path, the white-draw factors contributed to the numerator are
\[
3\cdot 4\cdots (k+2)=(3)_k,
\]
and the black-draw factors are
\[
2\cdot 3\cdots (N-k+1)=(2)_{N-k}.
\]
The denominator is always
\[
5\cdot 6\cdots (N+4)=(5)_N.
\]
Hence every such color word has probability
\[
\frac{(3)_k(2)_{N-k}}{(5)_N}.
\]

\subsection*{Step 2: sum over all words with $k$ white draws}
There are $\binom{N}{k}$ such color words. Therefore
\[
\mathbb{P}(I_N=k)=\binom{N}{k}\frac{(3)_k(2)_{N-k}}{(5)_N}.
\]
This is the beta-binomial distribution with parameters $(N,3,2)$.

Expanding the factorials gives the concrete formula
\begin{align*}
\mathbb{P}(I_N=k)
&=
\binom{N}{k}\frac{(k+2)!}{2!}\frac{(N-k+1)!}{1!}\frac{4!}{(N+4)!} \\
&=
\frac{12\,(k+1)(k+2)(N-k+1)}{(N+1)(N+2)(N+3)(N+4)}.
\end{align*}

\subsection*{Step 3: the finite-$N$ majority probability}
Thus
\[
\mathbb{P}\!\left(X_N>\frac12\right)
=
\sum_{k>\frac{N-1}{2}}
\frac{12\,(k+1)(k+2)(N-k+1)}{(N+1)(N+2)(N+3)(N+4)}.
\]

\section{Solution 2: A Riemann-sum asymptotic}

This route derives the limit directly from the exact finite law.

\subsection*{Step 1: recover the exact mass function locally}
As in Solution 1, any color word with exactly $k$ white draws and $N-k$ black draws has probability
\[
\frac{(3)_k(2)_{N-k}}{(5)_N},
\]
and there are $\binom{N}{k}$ such words. Hence
\[
\mathbb{P}(I_N=k)
=
\binom{N}{k}\frac{(3)_k(2)_{N-k}}{(5)_N}
=
\frac{12\,(k+1)(k+2)(N-k+1)}{(N+1)(N+2)(N+3)(N+4)}.
\]

\subsection*{Step 2: pass to a density}
Write $x=k/N$. Uniformly in $0\le k\le N$,
\[
\mathbb{P}(I_N=k)
=
\frac{12}{N}x^2(1-x)+O\!\left(\frac{1}{N^2}\right).
\]
Therefore
\[
\mathbb{P}\!\left(X_N>\frac12\right)
=
\sum_{k>\frac{N-1}{2}}
\left[
\frac{1}{N}\,12\left(\frac{k}{N}\right)^2\left(1-\frac{k}{N}\right)
+O\!\left(\frac{1}{N^2}\right)
\right].
\]
Since there are $O(N)$ summands, the total error is $O(1/N)$. Thus
\[
\mathbb{P}\!\left(X_N>\frac12\right)
\longrightarrow
\int_{1/2}^{1} 12x^2(1-x)\,dx.
\]

\subsection*{Step 3: evaluate the integral}
We compute
\[
12\int_{1/2}^{1}(x^2-x^3)\,dx
=
12\left[\frac{x^3}{3}-\frac{x^4}{4}\right]_{1/2}^{1}
=
\left[4x^3-3x^4\right]_{1/2}^{1}
=
1-\frac{5}{16}
=
\frac{11}{16}.
\]
Hence
\[
P_\infty=\frac{11}{16}.
\]

\section{Solution 3: Bayesian conjugacy and predictive equivalence}

This route interprets the urn as a mixture of Bernoulli sequences.

\subsection*{Step 1: the latent parameter}
Let $P\sim \mathrm{Beta}(3,2)$, with density
\[
f(p)=\frac{p^{3-1}(1-p)^{2-1}}{B(3,2)}=12p^2(1-p), \qquad 0<p<1.
\]
Conditional on $P=p$, let the colors be i.i.d.\ Bernoulli$(p)$, where ``white'' means success.

\subsection*{Step 2: posterior predictive probabilities}
After observing $k$ white draws in $N$ trials, the posterior is
\[
P\mid \text{data}\sim \mathrm{Beta}(3+k,\,2+N-k).
\]
Its mean is
\[
\mathbb{E}[P\mid \text{data}]
=
\frac{3+k}{5+N}.
\]
But this is exactly the physical probability of drawing a white ball from the urn when the current state is
\[
(W_N,B_N)=(3+k,\,2+N-k).
\]
Therefore the Bayesian predictive process and the P\'olya urn have the same finite-dimensional distributions.

\subsection*{Step 3: pass to the limit}
Conditional on $P$, the strong law of large numbers gives
\[
\frac{I_N}{N}\longrightarrow P
\qquad\text{almost surely.}
\]
Since
\[
X_N=\frac{3+I_N}{N+5}
\quad\text{and}\quad
\frac{I_N}{N}\to P,
\]
we also have
\[
X_N\longrightarrow P
\qquad\text{almost surely.}
\]
Thus
\[
X_\infty\sim \mathrm{Beta}(3,2),
\]
and
\[
P_\infty
=
\mathbb{P}\!\left(P>\frac12\right)
=
\int_{1/2}^{1}12p^2(1-p)\,dp
=
\frac{11}{16}.
\]

\section{Solution 4: Martingale convergence and factorial moments}

This proof identifies the limit law through martingales alone.

\subsection*{Step 1: $X_N$ is a bounded martingale}
Let $\mathcal{F}_N$ be the natural filtration. Since
\[
\mathbb{P}(W_{N+1}=W_N+1\mid \mathcal{F}_N)=\frac{W_N}{N+5},
\]
we have
\begin{align*}
\mathbb{E}[X_{N+1}\mid \mathcal{F}_N]
&=
\frac{W_N+1}{N+6}\cdot \frac{W_N}{N+5}
+
\frac{W_N}{N+6}\cdot \left(1-\frac{W_N}{N+5}\right) \\
&=
\frac{W_N}{N+5}
=
X_N.
\end{align*}
Thus $(X_N)$ is a martingale. Since $0\le X_N\le 1$, the martingale convergence theorem gives
\[
X_N\longrightarrow X_\infty
\qquad\text{almost surely and in }L^1.
\]

\subsection*{Step 2: factorial-moment martingales}
Fix $m\ge 1$ and define
\[
M_N^{(m)}:=\frac{(W_N)_m}{(N+5)_m}.
\]
We claim $(M_N^{(m)})$ is also a martingale.

Indeed, given $\mathcal{F}_N$, either $W_{N+1}=W_N+1$ with probability $W_N/(N+5)$, or $W_{N+1}=W_N$ with the complementary probability. Hence
\begin{align*}
\mathbb{E}\!\left[(W_{N+1})_m\mid \mathcal{F}_N\right]
&=
\frac{W_N}{N+5}(W_N+1)_m
+
\left(1-\frac{W_N}{N+5}\right)(W_N)_m.
\end{align*}
Using
\[
(W_N+1)_m=\frac{W_N+m}{W_N}(W_N)_m,
\]
we obtain
\begin{align*}
\mathbb{E}\!\left[(W_{N+1})_m\mid \mathcal{F}_N\right]
&=
\left(
\frac{W_N+m}{N+5}
+
1-\frac{W_N}{N+5}
\right)(W_N)_m \\
&=
\frac{N+5+m}{N+5}(W_N)_m.
\end{align*}
Therefore
\[
\mathbb{E}[M_{N+1}^{(m)}\mid \mathcal{F}_N]
=
\frac{1}{(N+6)_m}\cdot \frac{N+5+m}{N+5}(W_N)_m
=
\frac{(W_N)_m}{(N+5)_m}
=
M_N^{(m)}.
\]

\subsection*{Step 3: identify the moments of $X_\infty$}
Since $0\le M_N^{(m)}\le 1$ and
\[
\frac{(W_N)_m}{(N+5)_m}\longrightarrow X_\infty^m
\qquad\text{almost surely,}
\]
dominated convergence gives
\[
\mathbb{E}[X_\infty^m]
=
\lim_{N\to\infty}\mathbb{E}[M_N^{(m)}]
=
\mathbb{E}[M_0^{(m)}]
=
\frac{(3)_m}{(5)_m}.
\]
But these are exactly the moments of $\mathrm{Beta}(3,2)$:
\[
\mathbb{E}[P^m]=\frac{(3)_m}{(5)_m},
\qquad P\sim \mathrm{Beta}(3,2).
\]
Since probability measures on $[0,1]$ are uniquely determined by their moments (Hausdorff moment problem), it follows that
\[
X_\infty\sim \mathrm{Beta}(3,2).
\]
Therefore
\[
P_\infty=\mathbb{P}\!\left(X_\infty>\frac12\right)=\frac{11}{16}.
\]

\begin{remark}
There is also a different structural identification of the same Beta law via exchangeability and de Finetti's theorem. We present that separately in Solution 5.
\end{remark}

\section{Solution 5: Exchangeability and de Finetti's theorem}

This proof identifies the directing measure of the exchangeable color sequence.

\subsection*{Step 1: exchangeability from local path weights}
Fix a color word of length $N$ with exactly $k$ white draws and $N-k$ black draws. As in the exact path calculation,
\[
\mathbb{P}(\text{that word})
=
\frac{(3)_k(2)_{N-k}}{(5)_N}.
\]
This depends only on the counts $(k,N-k)$, not on the order of the colors. Therefore the infinite color sequence is exchangeable.

\subsection*{Step 2: de Finetti representation}
By de Finetti's theorem, there exists a unique probability measure $\mu$ on $[0,1]$ such that for all $N$ and $k$,
\[
\mathbb{P}(I_N=k)
=
\binom{N}{k}\int_0^1 p^k(1-p)^{N-k}\,d\mu(p).
\]
On the other hand, summing over the $\binom{N}{k}$ exchangeable words gives
\[
\mathbb{P}(I_N=k)
=
\binom{N}{k}\frac{(3)_k(2)_{N-k}}{(5)_N}.
\]
Hence
\[
\int_0^1 p^k(1-p)^{N-k}\,d\mu(p)
=
\frac{(3)_k(2)_{N-k}}{(5)_N}.
\]

\subsection*{Step 3: identify the directing measure}
Now
\begin{align*}
12\int_0^1 p^{k+2}(1-p)^{N-k+1}\,dp
&=
12\,B(k+3,N-k+2) \\
&=
12\,\frac{\Gamma(k+3)\Gamma(N-k+2)}{\Gamma(N+5)} \\
&=
\frac{(3)_k(2)_{N-k}}{(5)_N}.
\end{align*}
Therefore $\mu$ is exactly the Beta$(3,2)$ law:
\[
d\mu(p)=12p^2(1-p)\,dp.
\]

Conditional on the directing variable $P$, the colors are i.i.d.\ Bernoulli$(P)$, so the empirical fraction of white draws converges almost surely to $P$. Therefore $X_N\to P$ almost surely, and hence
\[
X_\infty\sim \mathrm{Beta}(3,2).
\]
Thus
\[
P_\infty=\frac{11}{16}.
\]

\section{Solution 6: Continuous-time embedding via Yule processes}

This route embeds the urn into continuous time.

\subsection*{Step 1: the continuous-time model}
Attach an independent rate-$1$ exponential clock to each ball. When a ball's clock rings, that ball produces one offspring of the same color. Let $W(t)$ and $B(t)$ be the white and black populations at time $t$.

At any moment, the total birth rate is $W(t)+B(t)$, and the probability that the next birth is white is
\[
\frac{W(t)}{W(t)+B(t)}.
\]
Therefore, if one observes the process only at birth times, the embedded jump chain is exactly the original P\'olya urn.

\subsection*{Step 2: Yule-process limits}
The white and black populations evolve independently as Yule processes started from $3$ and $2$ individuals. A standard theorem for Yule processes states that
\[
e^{-t}W(t)\longrightarrow \mathcal{W},
\qquad
e^{-t}B(t)\longrightarrow \mathcal{B}
\qquad\text{almost surely,}
\]
where
\[
\mathcal{W}\sim \mathrm{Gamma}(3,1),
\qquad
\mathcal{B}\sim \mathrm{Gamma}(2,1),
\]
and $\mathcal{W}$ and $\mathcal{B}$ are independent.

\subsection*{Step 3: pass to the ratio}
Let $\tau_N$ be the time of the $N$-th birth. Then $\tau_N\to\infty$ almost surely and
\[
W(\tau_N)=W_N,\qquad B(\tau_N)=B_N.
\]
Hence
\[
X_N
=
\frac{W_N}{W_N+B_N}
=
\frac{W(\tau_N)}{W(\tau_N)+B(\tau_N)}
\longrightarrow
\frac{\mathcal{W}}{\mathcal{W}+\mathcal{B}}
\qquad\text{almost surely.}
\]
By the gamma-beta algebra, if $U\sim \Gamma(\alpha,1)$ and $V\sim \Gamma(\beta,1)$ are independent, then
\[
\frac{U}{U+V}\sim \mathrm{Beta}(\alpha,\beta).
\]
Therefore
\[
X_\infty\sim \mathrm{Beta}(3,2),
\]
and
\[
P_\infty=\frac{11}{16}.
\]

\section{Solution 7: First tie, Catalan numbers, and a finite-time tie lemma}

This proof is self-contained and works directly with the finite-time probabilities
\[
\mathbb{P}\!\left(X_N>\frac12\right),
\]
without needing any continuity statement about the limiting law.

\subsection*{Step 1: majority probabilities from a tied start}

For $m\ge 1$, let $q_N^{(m)}$ denote the probability that, starting from a tied state $(m,m)$, the white proportion is strictly larger than $1/2$ after $N$ further draws.

\begin{lemma}[tie lemma]
For each fixed $m\ge 1$,
\[
q_N^{(m)}\longrightarrow \frac12
\qquad\text{as }N\to\infty.
\]
More precisely:
\begin{itemize}
    \item if $N=2n+1$ is odd, then $q_{2n+1}^{(m)}=\frac12$ exactly;
    \item if $N=2n$ is even, then
    \[
    q_{2n}^{(m)}
    =
    \frac12\left(1-\binom{2n}{n}\frac{(m)_n^2}{(2m)_{2n}}\right),
    \]
    and the tie probability
    \[
    \binom{2n}{n}\frac{(m)_n^2}{(2m)_{2n}}
    \longrightarrow 0.
    \]
\end{itemize}
\end{lemma}

\begin{proof}
Start from $(m,m)$.

If $N=2n+1$ is odd, then after $2n+1$ draws the difference between white and black counts is odd, so a tie is impossible. By color symmetry,
\[
q_{2n+1}^{(m)}=\frac12.
\]

If $N=2n$ is even, then white has strict majority exactly when there are more white draws than black draws among these $2n$ draws. By symmetry,
\[
q_{2n}^{(m)}
=
\frac12\Bigl(1-\mathbb{P}(\text{tie after }2n\text{ draws from }(m,m))\Bigr).
\]
A tie after $2n$ draws means exactly $n$ white draws and $n$ black draws. Every such color word has probability
\[
\frac{(m)_n(m)_n}{(2m)_{2n}},
\]
and there are $\binom{2n}{n}$ such words. Therefore
\[
\mathbb{P}(\text{tie after }2n\text{ draws from }(m,m))
=
\binom{2n}{n}\frac{(m)_n^2}{(2m)_{2n}}.
\]
Using Stirling's formula,
\[
\binom{2n}{n}\frac{(m)_n^2}{(2m)_{2n}}
=
\binom{2n}{n}
\frac{\Gamma(m+n)^2\Gamma(2m)}{\Gamma(m)^2\Gamma(2m+2n)}
=
O(n^{-1/2}),
\]
so it tends to $0$. Hence
\[
q_{2n}^{(m)}\longrightarrow \frac12.
\]
Combining the odd and even cases proves the lemma.
\end{proof}

\subsection*{Step 2: the first-tie distribution from the initial state $(3,2)$}

Now return to the original urn started from $(3,2)$, and let
\[
T:=\inf\{N\ge 0: W_N=B_N\}
\]
be the first time the process ever reaches a tie.

A first tie can only occur at some state
\[
(n+3,n+3),\qquad n\ge 0,
\]
after exactly $2n+1$ draws, consisting of
\[
n\ \text{white draws and}\ n+1\ \text{black draws}.
\]

To be a \emph{first} tie, the path must stay strictly above the diagonal until the final step. The number of such paths is the Catalan number
\[
C_n=\frac{1}{n+1}\binom{2n}{n}.
\]
Every such path has probability
\[
\frac{(3)_n(2)_{n+1}}{(5)_{2n+1}}
=
12\,\frac{(n+2)!^2}{(2n+5)!}.
\]
Therefore
\[
a_n:=\mathbb{P}\bigl(T=2n+1,\ (W_T,B_T)=(n+3,n+3)\bigr)
=
12\,C_n\,\frac{(n+2)!^2}{(2n+5)!}.
\]

Using
\[
C_n=\frac{(2n)!}{n!(n+1)!},
\]
this simplifies to
\[
a_n
=
\frac{3(n+2)}{(2n+1)(2n+3)(2n+5)}.
\]

\subsection*{Step 3: sum the first-tie probabilities}
A partial fraction decomposition gives
\[
\frac{3(n+2)}{(2n+1)(2n+3)(2n+5)}
=
\frac{9}{16(2n+1)}-\frac{3}{8(2n+3)}-\frac{3}{16(2n+5)}.
\]
Hence
\[
\sum_{n=0}^{\infty} a_n
=
\sum_{n=0}^{\infty}
\frac{3(n+2)}{(2n+1)(2n+3)(2n+5)}
=
\frac58.
\]
So
\[
\mathbb{P}(T<\infty)=\frac58.
\]

\subsection*{Step 4: decompose the finite-$N$ majority event}
For each $N\ge 0$,
\[
\mathbb{P}_{(3,2)}\!\left(X_N>\frac12\right)
=
\mathbb{P}(T>N)
+
\sum_{\substack{n\ge 0\\2n+1\le N}}
a_n\, q_{N-(2n+1)}^{(n+3)}.
\]
Indeed:
\begin{itemize}
    \item if $T>N$, then no tie has yet occurred, so the process is still strictly above the diagonal at time $N$, hence $X_N>1/2$;
    \item if the first tie occurs at time $2n+1$, then the process is at the tied state $(n+3,n+3)$, and the remaining probability of white majority after $N-(2n+1)$ further draws is exactly $q_{N-(2n+1)}^{(n+3)}$.
\end{itemize}

Now let $N\to\infty$. Since
\[
\mathbb{P}(T>N)\longrightarrow \mathbb{P}(T=\infty)=1-\sum_{n=0}^\infty a_n,
\]
and since each $q_{N-(2n+1)}^{(n+3)}\to \tfrac12$ by the tie lemma, dominated convergence yields
\begin{align*}
P_\infty
&=
\lim_{N\to\infty}\mathbb{P}_{(3,2)}\!\left(X_N>\frac12\right) \\
&=
\mathbb{P}(T=\infty)
+
\sum_{n=0}^\infty a_n\cdot \frac12 \\
&=
\left(1-\sum_{n=0}^\infty a_n\right)+\frac12\sum_{n=0}^\infty a_n \\
&=
1-\frac12\sum_{n=0}^\infty a_n.
\end{align*}
Since $\sum a_n=5/8$, we obtain
\[
P_\infty
=
1-\frac12\cdot\frac58
=
\frac{11}{16}.
\]

\section{Solution 8: Transport PDEs, raw weights, and Beta normalization}

This analytic route begins with an unnormalized generating function whose coefficients capture the correct state weights up to a factor depending only on the total level $N=i+j$. A Beta-kernel integral then restores the exact normalization.

\subsection*{Step 1: the unnormalized generating function}
For $i,j\ge 0$, define the raw weight
\[
u_{i,j}:=\frac{(3)_i(2)_j}{i!\,j!}
=
\binom{i+2}{2}(j+1).
\]
Consider the shifted generating function
\[
U(x,y,t):=\sum_{i,j\ge 0} u_{i,j}\,x^{i+3}y^{j+2}t^{i+j}.
\]
Summing the geometric series gives
\[
U(x,y,t)
=
x^3y^2\sum_{i\ge 0}\binom{i+2}{2}(xt)^i
\sum_{j\ge 0}(j+1)(yt)^j
=
\frac{x^3}{(1-xt)^3}\,\frac{y^2}{(1-yt)^2}.
\]

\subsection*{Step 2: the transport PDE}
The function $U$ satisfies the first-order PDE
\[
\frac{\partial U}{\partial t}
=
x^2\frac{\partial U}{\partial x}
+
y^2\frac{\partial U}{\partial y},
\qquad
U(x,y,0)=x^3y^2.
\]
Indeed, this follows either by direct differentiation of the closed form above or by the method of characteristics: along the characteristic curves
\[
\frac{dx}{dt}=-x^2,\qquad \frac{dy}{dt}=-y^2,
\]
the quantity $U$ remains constant, yielding exactly
\[
U(x,y,t)=\frac{x^3}{(1-xt)^3}\frac{y^2}{(1-yt)^2}.
\]

The coefficient of $x^{i+3}y^{j+2}t^{i+j}$ in $U$ is $u_{i,j}$, which captures the correct state dependence in $i$ and $j$ but is not yet normalized to a probability.

\subsection*{Step 3: identify the missing normalization}
The exact probability from Solution 1 is
\[
p_{i,j}:=\mathbb{P}(I_{i+j}=i)
=
\binom{i+j}{i}\frac{(3)_i(2)_j}{(5)_{i+j}}.
\]
Since
\[
\binom{i+j}{i}\frac{(3)_i(2)_j}{(5)_{i+j}}
=
u_{i,j}\cdot \frac{(i+j)!}{(5)_{i+j}},
\]
the discrepancy between $u_{i,j}$ and $p_{i,j}$ depends only on the total level
\[
N=i+j.
\]
For this urn,
\[
\frac{N!}{(5)_N}
=
\frac{24}{(N+1)(N+2)(N+3)(N+4)}
=
4\int_0^1 u^N(1-u)^3\,du.
\]

This suggests the integral operator
\[
(\mathcal T f)(t):=4\int_0^1 (1-u)^3 f(ut)\,du,
\]
which multiplies the coefficient of $t^N$ by the factor $N!/(5)_N$.

\subsection*{Step 4: Beta-kernel normalization}
Define the normalized shifted generating function
\[
\widetilde G(x,y,t):=(\mathcal T U)(x,y,t)
=
4\int_0^1 (1-u)^3 U(x,y,ut)\,du.
\]
Substituting the closed form for $U$ gives
\[
\widetilde G(x,y,t)
=
4x^3y^2
\int_0^1
\frac{(1-u)^3}{(1-uxt)^3(1-uyt)^2}\,du.
\]

Now extract coefficients under the integral:
\begin{align*}
[x^{i+3}y^{j+2}t^{i+j}]\,\widetilde G
&=
4\binom{i+2}{2}(j+1)\int_0^1 u^{i+j}(1-u)^3\,du \\
&=
4\binom{i+2}{2}(j+1)B(i+j+1,4) \\
&=
\frac{12(i+1)(i+2)(j+1)}{(i+j+1)(i+j+2)(i+j+3)(i+j+4)} \\
&=
p_{i,j}.
\end{align*}
Thus
\[
\widetilde G(x,y,t)
=
\sum_{i,j\ge 0} p_{i,j}\,x^{i+3}y^{j+2}t^{i+j}.
\]
In particular, the PDE route reproduces the exact beta-binomial law.

\subsection*{Step 5: recover the Beta-mixture form}
If one removes the harmless initial factor $x^3y^2$, the corresponding unshifted normalized generating function is
\[
G(x,y,t):=\sum_{i,j\ge 0} p_{i,j}\,x^i y^j t^{i+j}.
\]
Using the Beta-binomial formula and summing over $N=i+j$,
\begin{align*}
G(x,y,t)
&=
\sum_{N\ge 0}\sum_{i=0}^N
\binom{N}{i}
\left(\int_0^1 12\,p^{i+2}(1-p)^{N-i+1}\,dp\right)
x^i y^{N-i} t^N \\
&=
12\int_0^1 p^2(1-p)
\sum_{N\ge 0}\bigl(t(px+(1-p)y)\bigr)^N\,dp \\
&=
12\int_0^1 \frac{p^2(1-p)}{1-t(px+(1-p)y)}\,dp.
\end{align*}
This is exactly the generating-function form of the Beta$(3,2)$ mixture.

Therefore the coefficient array is beta-binomial, and the same limit law follows:
\[
X_\infty\sim \mathrm{Beta}(3,2).
\]
Hence
\[
P_\infty
=
\mathbb{P}\!\left(X_\infty>\frac12\right)
=
\int_{1/2}^1 12p^2(1-p)\,dp
=
\frac{11}{16}.
\]

\begin{remark}[optional normalized PDE]
One can also write down a first-order PDE directly for the normalized shifted generating function $\widetilde G$:
\[
t\,\widetilde G_t+4\widetilde G-4x^3y^2
=
t x^2 \widetilde G_x + t y^2 \widetilde G_y,
\qquad
\widetilde G(x,y,0)=x^3y^2.
\]
This equation is exact, but less transparent than the raw transport PDE for $U$. For this reason, the unnormalized PDE together with the Beta-kernel normalization is the cleaner analytic route.
\end{remark}

\section{Conclusion}

All eight routes lead to the same answer:
\[
\boxed{
\lim_{N\to\infty}\mathbb{P}\!\left(X_N>\frac12\right)=\frac{11}{16}
}.
\]

The common structural theme is the repeated appearance of the Beta$(3,2)$ law:
\begin{itemize}
    \item exactly, via the beta-binomial formula;
    \item asymptotically, via the density $12x^2(1-x)$;
    \item probabilistically, via Bayesian conjugacy, martingales, de Finetti, and Yule processes;
    \item analytically, via a transport PDE for raw weights and a Beta-kernel normalization.
\end{itemize}

\end{document}
